% ----- ----- ----- ----- ----- ----- ----- ----- ----- ----- 
% Preamble
% ----- ----- ----- ----- ----- ----- ----- ----- ----- ----- 
\documentclass[a4paper,9pt]{article}
\usepackage[utf8]{inputenc}
\usepackage[english]{babel}
\usepackage{amsmath}
\usepackage{authblk}
\usepackage{background}
\usepackage{enumitem}
\usepackage{fancyhdr}
\usepackage{geometry}
\usepackage{graphicx}
\usepackage[hidelinks]{hyperref}
\usepackage{indentfirst}
\usepackage{multicol}
\usepackage{tabularx}
\usepackage[most]{tcolorbox}
\usepackage{tgheros}
\usepackage[explicit]{titlesec}
\usepackage{xcolor}
% ----- ----- ----- ----- ----- 
% Column Spacing
% ----- ----- ----- ----- ----- 
\setlength{\columnseprule}{0pt}
\setlength{\columnsep}{10.0pt}
\newcolumntype{L}[1]{>{\raggedright\let\newline\\\arraybackslash\hspace{0pt}}m{#1}}
\newcolumntype{C}[1]{>{\centering\let\newline\\\arraybackslash\hspace{0pt}}m{#1}}
\newcolumntype{R}[1]{>{\raggedleft\let\newline\\\arraybackslash\hspace{0pt}}m{#1}}
% ----- ----- ----- ----- ----- 
% Custom Colors
% ----- ----- ----- ----- ----- 
\definecolor{myDColor}{HTML}{000000} % Darker Color
\definecolor{myLColor}{HTML}{CFB87C} % Lighter Color
\definecolor{LinkColor}{HTML}{CFB87C} % Link Color
% ----- ----- ----- ----- ----- 
% Custom Title
% ----- ----- ----- ----- ----- 
\makeatletter
\renewcommand*{\maketitle}{%
\noindent
\begin{minipage}{\textwidth}
\begin{tikzpicture}
\node[rectangle,rounded corners=10pt,inner sep=7.5pt,fill=myDColor,text width= 0.975\textwidth, align=center] {\color{white}\Huge \@title};
\end{tikzpicture}
\end{minipage}
\hfill
\bigskip\bigskip
}%
\makeatother
% ----- ----- ----- ----- ----- 
% Custom Section
% ----- ----- ----- ----- ----- 
\newcommand*\sectionlabel{}
\titleformat{\section}
{\gdef\sectionlabel{}
\normalfont\sffamily\Large\bfseries\scshape}
{\gdef\sectionlabel{\thesection\ }}{0pt}{
\noindent
\begin{tikzpicture}
\node[rectangle,rounded corners=6pt,inner sep=4pt,fill=myLColor,text width= 0.975\columnwidth, align=center] {\color{white}\sectionlabel#1};
\end{tikzpicture}}
\titlespacing*{\section}{0pt}{5pt}{5pt}
% ----- ----- ----- ----- ----- 
% Custom Header & Footer / Geometry
% ----- ----- ----- ----- ----- 
\geometry{a4paper,left=10mm,right=10mm,top=15mm,bottom=15mm}
\linespread{1} % Distance between lines
\pagestyle{fancy}
\fancyhead[L]{Taylor Larrechea}
\fancyhead[C]{Linked List Assignment Interview}
\fancyhead[R]{University of Colorado}
\fancyfoot[L]{CSPB 2270 Data Structures}
\fancyfoot[R]{Linked List AI}
\backgroundsetup{
    scale = 1,
    angle = 0,
    opacity = 0.1,
    contents = {
    \includegraphics[scale = 0.25, keepaspectratio]{Figures/CU Seal.png}
    }
}
\newtcolorbox{highlight}[1][]{%
    enhanced,
    skin first=enhanced,
    skin middle=enhanced,
    skin last=enhanced,
    before upper={\parindent15pt},
    breakable,
    boxrule = 0pt,
    frame hidden,
    borderline west = {4pt}{0pt}{myDColor},
    colback = myLColor!10,
    coltitle = myLColor!5,
    sharp corners,
    rounded corners = southeast,
    arc is angular,
    arc = 3mm,
    attach boxed title to top left,
    boxed title style = {%
        enhanced,
        colback = myDColor,
        colframe = myDColor,
        top = 0pt,
        bottom = 0pt,
        sharp corners,
        rounded corners = northeast,
        arc is angular,
        arc = 2mm,
        rightrule = 0pt,
        bottomrule = 0pt,
        toprule = 0pt,
    },
    title = {\bfseries\large #1}, 
    overlay unbroken={%
        \path[fill = tcbcolback!80!black] 
            ([yshift = 3mm]interior.south east) -- ++(-0.4,-0.1) -- ++(0.1,-0.2);
    },
    overlay first = {%
        \path[fill = tcbcolback!80!black] 
            ([yshift = 3mm]interior.south east) -- ++(-0.4,-0.1) -- ++(0.1,-0.2);
    },
    overlay middle={%
        \path[fill = tcbcolback!80!black] 
            ([yshift = -3mm]interior.north east) -- ++(-0.4,0.1) -- ++(0.1,0.2);
        \path[fill = tcbcolback!80!black] 
            ([yshift = 3mm]interior.south east) -- ++(-0.4,-0.1) -- ++(0.1,-0.2);
    },
    overlay last={%
        \path[fill = tcbcolback!80!black] 
            ([yshift = -3mm]interior.north east) -- ++(-0.4,0.1) -- ++(0.1,0.2);
        \path[fill = tcbcolback!80!black] 
            ([yshift = 3mm]interior.south east) -- ++(-0.4,-0.1) -- ++(0.1,-0.2);
    },
    extras middle and last = { rounded corners = northeast }
}
% ----- ----- ----- ----- ----- Title
\title{Linked List Assignment Interview Notes}
% ----- ----- ----- ----- ----- Begin Document
\begin{document}

\maketitle

\section*{Vector Overview}

\subsection*{Generic Overview}

A vector is a dynamic array that allows storing and accessing elements of the same data type. It is a container class
provided by the Standard Template Library (STL) in C++. Vectors provide several advantages over traditional arrays, such
as automatic memory management, resizable capacity, and built-in functions for efficient manipulation.

Internally, a vector is implemented as a contiguous block of memory that stores the elements. It dynamically manages its
size, expanding or shrinking as needed. Elements are accessed using zero-based indexing, similar to arrays. The vector
class provides various member functions and operators for performing common operations like adding elements at the end,
inserting elements at specific positions, removing elements, and accessing elements by index. Vectors offer efficient
random access, constant time complexity for accessing elements, and amortized constant time complexity for adding or 
removing elements at the end.

\subsection*{How Are Vectors Different From Arrays?}

\noindent Vectors and arrays are both used to store collections of elements in C++, but vary in some areas:

\begin{itemize}
  \item Arrays are fixed-size collections with a static size determined at compile-time, while vectors are dynamic arrays
  with a flexible and resizable size.
  \item Arrays require manual memory management and cannot be easily resized, whereas vectors handle memory allocation and
  deallocation automatically, allowing for dynamic resizing.
  \item Arrays provide direct and efficient access to elements using zero-based indexing, while vectors offer convenient
  member functions like `push\_back()' and `pop\_back()' for adding and removing elements.
  \item Arrays have a fixed memory layout, resulting in slightly faster element access, while vectors provide more convenience
  and flexibility.
  \item Arrays cannot be easily copied or assigned to another array, while vectors can be easily copied of assigned to other
  vectors.
\end{itemize}

\noindent In summary, arrays are suitable when the size is known and fixed, and direct memory control is needed. Vectors, on the other hand,
are preferred when the size needs to be dynamic, convenient operations are required, and automatic memory management is desired.

\subsection*{What Is The Algorithm For Adding And Removing Elements From A Vector?}

\noindent To add elements to a vector in C++, we use:

\begin{itemize}
  \item To add an element at the end of the vector, we use the `push\_back()' function, which appends the element to the end of
  the vector.
  \item To insert an element at a specific position within the vector, we use `insert()' function. This function takes an iterator
  pointing to the position and the element to be inserted.
\end{itemize}

\noindent To remove elements from a vector in C++, we use:

\begin{itemize}
  \item To remove the last element of the vector, we use the `pop\_back()' function, which removes the element from the end of the vector.
  \item If you want to remove an element from a specific position within the vector, we use the `erase()' function. It takes an iterator
  pointing to the position and removes the element at that position.
\end{itemize}

\noindent These operations provide a convenient and efficient way to add and remove elements from a vector, maintaining the integrity and
resizing the vector dynamically as needed.

\subsection*{What Are The Benefits Of Using A Vector Over An Array?}

\noindent The benefits of using a vector over an array are:

\begin{itemize}
  \item Dynamic Size: Vectors provide dynamic resizing, allowing you to add or remove elements at runtime. Unlike arrays, which have a fixed
  size, vectors automatically handle memory allocation.
  \item Convenience: Vectors come with built-in functions like `push\_back()', `pop\_back()', and `insert()', making it easier to add, remove,
  or insert elements without manual management. Vectors also provide member functions like `size()' and `empty()' to retrieve information about
  the number of elements and check if the vector is empty.
  \item Bounds Checking: Vectors perform bounds checking on element access, ensuring that you do not go out of bounds. This helps prevent accessing
  memory beyond the allocated size and reduces the risk of errors and crashes that can occur with arrays.
  \item Iteration Support: Vectors can be easily iterated using range-based for loops or iterators, providing a convenient way to traverse and
  manipulate elements.
  \item Copying and Assignment: Vectors can be easily copied and assigned to other vectors using the copy constructor and assignment operator, allowing
  for efficient and convenient data handling.
\end{itemize}

\noindent Overall, vectors offer the flexibility of dynamic resizing, convenient functions, bounds checking, iteration support, and ease of copying,
making them a preferred choice over arrays in many scenarios.

\subsection*{What Is The Basic Algorithm For Iterating Through A Vector?}

The simple algorithm for iterating through the elements of a vector is to use a for-loop. Starting the index at 0, and going up to the size of
the vector, we can use built in functions like `at()' to indicate which index we are wanting to examine / manipulate. Using this simple algorithm, we check
for certain conditions within our vector, look for specific elements, and many other things as well.

\section*{Linked List Overview}

\subsection*{Generic Overview}

A linked list is a data structure commonly used in object-oriented programming that consists of nodes connected through 
pointers. Each node contains a value and a pointer to the next node in the list, except for the last node which points 
to `nullptr'. The linked list does not have a specific `head' or `tail' node, but rather, the first node in the list is 
often referred to as the `head' and the last node is commonly known as the `tail'. The nodes in between the `head' and 
`tail' are linked in a sequential manner. The `head' node serves as the starting point for traversing the list, and each 
node points to the next node, allowing efficient insertion and removal operations.

\end{document}
% ----- ----- ----- ----- ----- End Document